\documentclass{exam}
\usepackage{ModeloIFBA}
%\usepackage[alf]{abntex2cite}
\usepackage{multicol}
\usepackage{adjustbox}
\usepackage{multicol}

\setlength\columnsep{30pt}


\renewcommand{\curso}{Tecnologia em Análise e Desenvolvimento de Sistemas }
\renewcommand{\materia}{Programação Orientada a Objetos}
\renewcommand{\turma}{2025.1}
\renewcommand{\professor}{Leandro Costa Souza}
\renewcommand{\avaliacao}{Prova Unidade II}
\renewcommand{\dataavaliacao}{01/09/2025}


\newcommand{\duracao}{1h 40min}
\newcounter{solcount}
\setcounter{solcount}{1}
\newcommand\CC{\
    \CorrectChoice\label{sol:\arabic{solcount}}
    \stepcounter{solcount}
}

\newcommand\printmyanswers{
\newpage
\section*{\centering Answers Sheet}
    \foreach \x in {1,...,\thequestion}{
    \noindent
        \x.\ref{sol:\x}\\par
    }
}

\printanswers

\begin{document}
\begin{adjustbox}{max width=\textwidth}
\begin{tabularx}{\textwidth}{rrXrl}
\multirow{5}{*}{\ifbalogo{1.7in}} & \textbf{Curso:} & \multicolumn{3}{l}{ \curso } \\
& \textbf{Matéria:} & \materia \space & \textbf{Turma:} & \small{\turma}  \\
& \textbf{Professor:} & \professor & & \\
& \textbf{Aluno:} & \hrulefill & \textbf{Data:} & \dataavaliacao \\
\end{tabularx}
\end{adjustbox}
\begin{tabularx}{\textwidth}{XcX}
 & \textbf{\avaliacao} &  \\ \hline
\end{tabularx}	
	\input{orientacoes}

	\begin{questions}
\begin{multicols}{2}

% --- QUESTÕES EXISTENTES (1-100) ---

\question[1] O que a sigla JDK significa no contexto do desenvolvimento Java?
\begin{choices}
 \choice Java Runtime Kit - Kit de Execução Java
 \CC Java Development Kit - Kit de Desenvolvimento Java
 \choice Java Deployment Kit - Kit de Implantação Java
 \choice Java Dynamic Kit - Kit Dinâmico Java
 \choice Java Developer Kernel - Núcleo de Desenvolvimento Java
\end{choices}

\question[1] Qual é a função da JVM (Java Virtual Machine)?
\begin{choices}
 \choice Compilar o código-fonte Java (.java) em bytecode (.class).
 \CC Fornecer um ambiente para executar o bytecode Java em qualquer plataforma.
 \choice Depurar o código Java durante o desenvolvimento.
 \choice Gerenciar as dependências do projeto.
 \choice Escrever código Java automaticamente.
\end{choices}

\question[1] Qual dos seguintes tipos de dados em Java é usado para armazenar números inteiros?
\begin{choices}
 \choice double
 \choice float
 \CC int
 \choice char
 \choice boolean
\end{choices}

\question[1] Qual estrutura de repetição é mais adequada quando não se sabe o número exato de iterações, mas se conhece a condição de parada?
\begin{choices}
 \choice for
 \choice foreach
 \CC while
 \choice if-else
 \choice switch
\end{choices}


\question[1] Qual é a convenção de nomenclatura para variáveis e métodos em Java?
\begin{choices}
 \choice PascalCase (Ex: MinhaVariavel)
 \choice snake\_case (Ex: minha\_variavel)
 \choice kebab-case (Ex: minha-variavel)
 \CC camelCase (Ex: minhaVariavel)
 \choice UPPERCASE (Ex: MINHAVARIAVEL)
\end{choices}

\question[1] O que é o escopo de uma variável?
\begin{choices}
 \choice O valor máximo que a variável pode armazenar.
 \choice O tipo de dado da variável.
 \CC A região do código onde a variável é acessível.
 \choice O nome único da variável.
 \choice O tempo de vida da variável na memória.
\end{choices}


\question[1] O que é uma classe em POO?
\begin{choices}
 \choice Um objeto em execução.
 \CC Um modelo para criar objetos.
 \choice Uma função que retorna um valor.
 \choice Uma variável que armazena dados.
 \choice Um evento do sistema.
\end{choices}

\question[1] O que são atributos em uma classe?
\begin{choices}
 \choice As ações que um objeto pode realizar.
 \CC As características ou dados que um objeto possui.
 \choice Os eventos que um objeto pode disparar.
 \choice Os construtores da classe.
 \choice As classes das quais ela herda.
\end{choices}

\question[1] Qual é o objetivo de um construtor em uma classe Java?
\begin{choices}
 \choice Destruir um objeto quando ele não é mais necessário.
 \CC Inicializar um objeto quando ele é criado.
 \choice Definir os métodos que o objeto pode executar.
 \choice Retornar o valor de um atributo privado.
 \choice Comparar dois objetos.
\end{choices}

\question[1] O que é encapsulamento?
\begin{choices}
 \choice A capacidade de uma classe herdar características de outra.
 \CC A prática de agrupar dados e métodos em uma unidade, escondendo os detalhes internos.
 \choice A capacidade de um objeto ter múltiplas formas.
 \choice A criação de múltiplos objetos a partir da mesma classe.
 \choice A organização de classes em pacotes.
\end{choices}

\question[1] Qual modificador de acesso torna um atributo ou método acessível apenas instâncias da própria classe?
\begin{choices}
 \choice public
 \choice protected
 \choice default
 \CC private
 \choice static
\end{choices}

\question[1] Para que servem os métodos `get' e `set' (getters e setters)?
\begin{choices}
 \choice Para criar e destruir objetos, respectivamente.
 \CC Para fornecer acesso controlado aos atributos privados de uma classe.
 \choice Para comparar e definir a ordem dos objetos.
 \choice Para iniciar e parar a execução de threads.
 \choice Para converter objetos em Strings.
\end{choices}

\question[1] O que é herança em POO?
\begin{choices}
 \choice A criação de uma cópia exata de um objeto.
 \CC Um mecanismo que permite a uma classe adquirir os atributos e métodos de outra.
 \choice A capacidade de um método ter o mesmo nome mas diferentes parâmetros.
 \choice A prática de esconder a complexidade interna de um objeto.
 \choice A relação onde um objeto contém outro.
\end{choices}

\question[1] Qual palavra-chave é usada em Java para indicar que uma classe herda de outra?
\begin{choices}
 \choice implements
 \choice inherits
 \CC extends
 \choice super
 \choice derives
\end{choices}

\question[1] O que a anotação `@Override' indica?
\begin{choices}
 \choice Que o método está sobrecarregado.
 \CC Que o método está sendo sobrescrito de uma superclasse.
 \choice Que o método é a versão final e não pode ser modificado.
 \choice Que o método é obsoleto e não deve ser usado.
 \choice Que o método é estático.
\end{choices}

\question[1] Qual é a saída do código a seguir?
\lstinputlisting{code/Veiculo.java}
\begin{choices}
 \choice `Veículo se move'
 \CC `Carro acelera'
 \choice Erro de compilação
 \choice `Veículo se move'`Carro acelera'
 \choice Nenhuma saída
\end{choices}

\question[1] Qual a diferença entre sobrecarga (overloading) e sobrescrita (overriding)?
\begin{choices}
 \choice Sobrecarga é em classes diferentes, sobrescrita é na mesma classe.
 \choice Sobrecarga usa a mesma assinatura, sobrescrita usa assinaturas diferentes.
 \CC Sobrecarga é ter métodos com mesmo nome e parâmetros diferentes na mesma classe; sobrescrita é redefinir um método da superclasse na subclasse.
 \choice Sobrecarga é para construtores, sobrescrita é para métodos comuns.
 \choice Não há diferença, são termos sinônimos.
\end{choices}

\question[1] O que a palavra-chave `super' faz dentro de um método de uma subclasse?
\begin{choices}
 \choice Chama um método da própria classe (subclasse).
 \choice Refere-se à instância atual do objeto.
 \CC Chama um método ou construtor da superclasse.
 \choice Cria uma nova instância da superclasse.
 \choice Finaliza a execução do método.
\end{choices}

\question[1] Em UML, qual tipo de seta representa a relação de Herança?
\begin{choices}
 \choice Seta pontilhada com ponta aberta.
 \choice Linha contínua com um losango vazio na ponta.
 \choice Linha contínua com um losango preenchido na ponta.
 \CC Linha contínua com uma seta triangular fechada e vazia.
 \choice Linha contínua com uma seta simples na ponta.
\end{choices}

\question[1] Qual a diferença fundamental entre Agregação e Composição?
\begin{choices}
 \choice Na Composição, as partes podem existir sem o todo; na Agregação, não.
 \CC Na Agregação, as partes podem existir sem o todo; na Composição, as partes dependem do ciclo de vida do todo.
 \choice Agregação é uma relação ``é um'', enquanto Composição é ``tem um''.
 \choice Não há diferença, são a mesma coisa.
 \choice Composição usa a palavra-chave `composes', enquanto Agregação usa `aggregates'.
\end{choices}


\question[1] O que o método `toString()' da classe `Object' faz por padrão?
\begin{choices}
 \choice Retorna uma representação em String vazia do objeto.
 \CC Retorna o nome da classe seguido por "@" e o código hash do objeto.
 \choice Retorna `null'.
 \choice Lança uma exceção.
 \choice Retorna uma lista de todos os atributos e seus valores.
\end{choices}

\question[1] Se uma classe é declarada como `final', o que isso significa?
\begin{choices}
 \choice Que ela não pode ser instanciada.
 \choice Que todos os seus métodos são estáticos.
 \CC Que ela não pode ser herdada por outra classe.
 \choice Que ela só pode ter um único objeto.
 \choice Que seus atributos não podem ser modificados.
\end{choices}

\question[1] Qual é o resultado de comparar duas referências de objetos com `=='?
\begin{choices}
 \choice Verifica se os objetos são do mesmo tipo.
 \choice Verifica se os valores dos atributos dos objetos são iguais.
 \CC Verifica se as duas referências apontam para o mesmo objeto na memória.
 \choice Chama o método `.equals()' automaticamente.
 \choice Sempre retorna `false'.
\end{choices}

\question[1] Para comparar o conteúdo de dois objetos do tipo `String', qual método deve ser usado?
\begin{choices}
 \choice `=='
 \choice `compare()'
 \choice `isEqual()'
 \CC `equals()'
 \choice `match()'
\end{choices}

\question[1] No Greenfoot, qual classe representa o cenário do jogo?
\begin{choices}
 \choice `Actor'
 \choice `Stage'
 \choice `Scene'
 \CC `World'
 \choice `Canvas'
\end{choices}

\question[1] Para criar um objeto que pode ser adicionado ao mundo, sua classe deve herdar de:
\begin{choices}
 \choice `World'
 \choice `Object'
 \CC `Actor'
 \choice `Sprite'
 \choice `Entity'
\end{choices}

\question[1] Qual método é chamado repetidamente em cada ator, constituindo o loop principal do jogo?
\begin{choices}
 \choice `run()'
 \choice `update()'
 \choice `loop()'
 \CC `act()'
 \choice `execute()'
\end{choices}

\question[1] Qual método do Greenfoot é usado para verificar se uma tecla específica está sendo pressionada?
\begin{choices}
 \choice `Greenfoot.isKeyPressed()'
 \CC `Greenfoot.isKeyDown()'
 \choice `Greenfoot.getKey()'
 \choice `Greenfoot.checkKey()'
 \choice `Greenfoot.isPressed()'
\end{choices}

\question[1] Qual método é usado para detectar se um ator está tocando em outro de uma classe específica?
\begin{choices}
 \choice `isColliding(Classe.class)'
 \choice `intersects(Classe.class)'
 \choice `hasContact(Classe.class)'
 \choice `onCollision(Classe.class)'
 \CC `isTouching(Classe.class)'
\end{choices}

\question[1] Como você removeria do mundo um objeto da classe `Moeda` que o ator atual está tocando?
\begin{choices}
 \choice `delete(Moeda.class);'
 \CC `removeTouching(Moeda.class);'
 \choice `getWorld().removeObject( getTouching( Moeda.class));'
 \choice `destroy(Moeda.class);'
 \choice `getWorld().removeTouching(Moeda.class);'
\end{choices}

\question[1] Para adicionar um novo ator (ex: `new Bala()`) no mundo na posição (100, 150) a partir de um ator, qual comando é usado?
\begin{choices}
 \choice `addObject(new Bala(), 100, 150);'
 \CC `getWorld().addObject(new Bala(), 100, 150);'
 \choice `Greenfoot.addObject(new Bala(), 100, 150);'
 \choice `new Bala(100, 150);'
 \choice `World.add(new Bala(), 100, 150);'
\end{choices}

\question[1] Qual método do Greenfoot gera um número inteiro aleatório até um limite (exclusive)?
\begin{choices}
 \choice `Greenfoot.random(limit);'
 \choice `Greenfoot.getRandom(limit);'
 \CC `Greenfoot.getRandomNumber(limit);'
 \choice `Greenfoot.nextInt(limit);'
 \choice `Greenfoot.randomInteger(limit);'
\end{choices}

\question[1] Se você tem várias classes de inimigos que compartilham comportamentos, qual seria a melhor abordagem de POO?
\begin{choices}
 \choice Copiar e colar o código comum em todas as classes.
 \CC Criar uma superclasse `Inimigo' com o comportamento comum e fazer as classes de inimigos herdarem dela.
 \choice Criar um método estático em uma classe `Util' para o comportamento comum.
 \choice Colocar todo o código na classe `World'.
 \choice Usar uma interface para cada tipo de comportamento.
\end{choices}

\question[1] Para tocar um arquivo de som chamado ``explosao.wav'', qual comando você usaria?
\begin{choices}
 \choice `Greenfoot.play("explosao.wav");'
 \CC `Greenfoot.playSound("explosao.wav");'
 \choice `Greenfoot.sound("explosao.wav");'
 \choice `Greenfoot.startSound("explosao.wav");'
 \choice `Greenfoot.music("explosao.wav");'
\end{choices}

\question[1] O que o método `getOneIntersectingObject(Classe.class)' retorna se não há intersecção?
\begin{choices}
 \choice `false`
 \choice `0`
 \CC `null`
 \choice Lança uma exceção `NullPointerException`.
 \choice `-1`
\end{choices}

\question[1] Como você pode obter a largura do mundo atual dentro de um ator?
\begin{choices}
 \choice `getWorld().width()'
 \choice `World.getWidth()'
 \CC `getWorld().getWidth()'
 \choice `Greenfoot.getWorldWidth()'
 \choice `this.World.width'
\end{choices}

\question[1] No tutorial do jogo "Hole.io", por que foi criada a classe `Swallowable'?
\begin{choices}
 \choice Para representar o jogador.
 \CC Para ser uma superclasse para todos os objetos que podem ser "engolidos", aplicando herança e polimorfismo.
 \choice Para controlar o placar do jogo.
 \choice Para definir os limites do mundo.
 \choice Para tocar os sons do jogo.
\end{choices}

\question[1] Para fazer um ator se mover em direção ao mouse, qual informação você precisaria obter primeiro?
\begin{choices}
 \choice A posição do centro do mundo.
 \choice O número de atores no cenário.
 \CC As informações do mouse, usando `Greenfoot.getMouseInfo()'.
 \choice A última tecla pressionada pelo usuário.
 \choice A velocidade de execução do cenário.
\end{choices}

\question[1] Qual a principal diferença entre tipos primitivos e tipos por referência em Java?
\begin{choices}
 \choice Tipos primitivos são mais rápidos.
 \choice Tipos primitivos armazenam o valor, enquanto tipos por referência armazenam um endereço de memória.
 \choice Tipos primitivos não podem ter métodos.
 \choice Tipos por referência são sempre mutáveis.
 \CC As alternativas 'b' e 'c' estão corretas.
\end{choices}

\question[1] O que o JRE (Java Runtime Environment) inclui?
\begin{choices}
 \choice Apenas o compilador Java (javac).
 \CC A JVM e as bibliotecas padrão do Java, mas não as ferramentas de desenvolvimento.
 \choice O JDK completo.
 \choice Apenas a JVM.
 \choice Ferramentas como o VSCode.
\end{choices}

\question[1] Qual é a saída do código a seguir?
\lstinputlisting{code/Ola.java}
\begin{choices}
 \choice true
 \CC false
 \choice Erro de compilação
 \choice `Olá'
 \choice Nada é impresso.
\end{choices}

\question[1] Qual operador é usado para negação lógica em Java?
\begin{choices}
 \choice `\&\&'
 \choice `||'
 \CC `!'
 \choice `\~'
 \choice `NOT'
\end{choices}


\question[1] Como se declara uma constante em Java?
\begin{choices}
 \choice Usando a palavra-chave `const`.
 \CC Usando as palavras-chave `final` e `static`.
 \choice Apenas declarando a variável com letras maiúsculas.
 \choice Usando a palavra-chave `constant`.
 \choice Não é possível declarar constantes em Java.
\end{choices}


\question[1] O que é a "pilha de execução" (stack) em Java?
\begin{choices}
 \choice Uma estrutura de dados para armazenar objetos criados com `new`.
 \CC Uma área de memória onde são armazenadas as chamadas de métodos e variáveis locais.
 \choice Uma fila de processos esperando para serem executados pela JVM.
 \choice O local onde o bytecode é armazenado.
 \choice Uma biblioteca para operações matemáticas.
\end{choices}

\question[1] O que é um objeto em POO?
\begin{choices}
 \choice Um modelo para criar outras classes.
 \CC Uma instância de uma classe.
 \choice Um tipo de dado primitivo.
 \choice Um método estático.
 \choice Um pacote de classes.
\end{choices}

\question[1] O que acontece se uma classe não define explicitamente um construtor?
\begin{choices}
 \choice A classe não pode ser instanciada.
 \choice Ocorre um erro de compilação.
 \CC O compilador Java fornece um construtor padrão (sem argumentos) automaticamente.
 \choice É necessário chamar um método `init()` para criar o objeto.
 \choice A classe só pode ter métodos estáticos.
\end{choices}

\question[1] O que a palavra-chave `this' representa dentro de um método de instância?
\begin{choices}
 \choice A superclasse.
 \choice A classe atual.
 \CC A instância atual do objeto sobre o qual o método foi chamado.
 \choice Um novo objeto que será criado.
 \choice O pacote onde a classe está localizada.
\end{choices}

\question[1] Uma classe `Carro' "tem um" `Motor'. Que tipo de relação de POO isso representa?
\begin{choices}
 \choice Herança
 \choice Polimorfismo
 \choice Abstração
 \CC Associação (Agregação ou Composição)
 \choice Interface
\end{choices}

\question[1] Qual é a classe raiz (superclasse) de todas as classes em Java?
\begin{choices}
 \choice `Main'
 \choice `System'
 \choice `Class'
 \CC `Object'
 \choice `Java'
\end{choices}


\question[1] Qual é a saída do código a seguir?
\lstinputlisting{code/Pai.java}
\begin{choices}
 \choice `Filho'
 \choice `Pai'
 \choice `Filho' `Pai'
 \CC `Pai' `Filho'
 \choice Erro de compilação
\end{choices}

\question[1] Se um método na superclasse é declarado como `private', ele pode ser sobrescrito na subclasse?
\begin{choices}
 \choice Sim, mas ele também deve ser `private' na subclasse.
 \choice Sim, e pode ter qualquer modificador de acesso na subclasse.
 \CC Não, porque o método não é visível para a subclasse.
 \choice Apenas se a subclasse estiver no mesmo pacote.
 \choice Sim, mas causa um aviso de compilação.
\end{choices}

\question[1] O que é "acoplamento" em POO?
\begin{choices}
 \choice O grau de responsabilidade de uma única classe.
 \choice A forma como os objetos são ligados na memória.
 \CC O grau de interdependência entre as classes. Baixo acoplamento é desejável.
 \choice O processo de juntar classes em um pacote.
 \choice A velocidade com que os métodos de uma classe são executados.
\end{choices}

\question[1] O que é "coesão" em POO?
\begin{choices}
 \choice O grau de interdependência entre as classes.
 \CC A medida em que os elementos de uma classe são focados em uma única responsabilidade.
 \choice A relação entre uma superclasse e uma subclasse.
 \choice A prática de esconder os dados de uma classe.
 \choice A capacidade de um objeto assumir várias formas.
\end{choices}

\question[1] Qual a finalidade da palavra reservada `package' em Java?
\begin{choices}
 \choice Definir o ponto de entrada de um programa.
 \choice Criar um arquivo executável.
 \choice Declarar variáveis globais.
 \choice Controlar o fluxo de execução do programa.
 \CC Agrupar classes relacionadas e evitar conflitos de nomes.
\end{choices}

\question[1] Se um atributo é declarado como `protected', onde ele pode ser acessado?
\begin{choices}
 \choice Apenas dentro da própria classe.
 \choice Em qualquer lugar do programa.
 \CC Dentro da própria classe, por subclasses e por outras classes no mesmo pacote.
 \choice Apenas por subclasses.
 \choice Apenas por outras classes no mesmo pacote.
\end{choices}


\question[1] Qual é o propósito do construtor de uma classe `World' no Greenfoot?
\begin{choices}
 \choice Definir a imagem de fundo do ator.
 \CC Inicializar o tamanho do mundo (largura, altura) e o tamanho da célula.
 \choice Adicionar os atores ao mundo.
 \choice Iniciar o loop do jogo.
 \choice Definir a velocidade de execução do cenário.
\end{choices}

\question[1] O que o método `turn(int amount)' faz em um ator?
\begin{choices}
 \choice Move o ator para a frente na quantidade especificada.
 \choice Vira o ator para um ângulo absoluto.
 \CC Gira a imagem do ator em um determinado número de graus a partir de sua orientação atual.
 \choice Define a transparência da imagem do ator.
 \choice Vira a imagem do ator horizontalmente.
\end{choices}

\question[1] Qual método retorna uma referência ao mundo em que o ator está?
\begin{choices}
 \choice `getParent()'
 \choice `getScene()'
 \choice `getContext()'
 \CC `getWorld()'
 \choice `getEnvironment()'
\end{choices}

\question[1] O que a classe `MouseInfo' do Greenfoot permite obter?
\begin{choices}
 \choice Apenas se o botão do mouse foi clicado.
 \CC As coordenadas X e Y do ponteiro do mouse, qual botão foi pressionado e qual ator está sob o mouse.
 \choice A marca e o modelo do mouse do usuário.
 \choice A velocidade de movimento do mouse.
 \choice O número de cliques realizados durante o jogo.
\end{choices}

\question[1] Se você quer que um ator faça algo apenas uma vez quando é adicionado ao mundo, onde você colocaria esse código?
\begin{choices}
 \choice No método `act()`, dentro de um `if`.
 \CC No construtor da classe do ator.
 \choice No construtor da classe `World`.
 \choice Na classe pai `Actor'.
\end{choices}

\question[1] Como você pode remover o próprio ator de dentro de sua classe?
\begin{choices}
 \choice `remove(this);`
 \choice `delete this;`
 \CC `getWorld().removeObject(this);`
 \choice `this.destroy();`
 \choice `Greenfoot.removeActor(this);`
\end{choices}

\question[1] Qual é o sistema de coordenadas do Greenfoot?
\begin{choices}
 \choice O ponto (0,0) está no centro do mundo.
 \choice O ponto (0,0) está no canto inferior esquerdo, com Y para cima.
 \choice O sistema de coordenadas é polar.
 \choice O sistema de coordenadas é definido pelo programador.
 \CC O ponto (0,0) está no canto superior esquerdo, com X para a direita e Y para baixo.
\end{choices}

% --- NOVAS QUESTÕES (101-150) ---


\question[1] No contexto da JVM, JRE e JDK, qual deles é necessário apenas para executar uma aplicação Java, sem a necessidade de compilá-la?
\begin{choices}
 \choice JVM
 \CC JRE
 \choice JDK
 \choice Javac
 \choice Maven
\end{choices}

\question[1] Qual é a saída do código a seguir?
\lstinputlisting{code/For.java}
\begin{choices}
 \choice 4
 \CC 5
 \choice 0
 \choice Erro de compilação, `i' fora de escopo.
 \choice 6
\end{choices}

\question[1] Qual a principal característica de um método `final'?
\begin{choices}
 \choice Ele não pode ser chamado.
 \CC Ele não pode ser sobrescrito (`overridden') por subclasses.
 \choice Ele só pode ser chamado uma vez.
 \choice Ele deve ser estático.
 \choice Ele não pode ter parâmetros.
\end{choices}


\question[1] Como você pode fazer um ator se mover para trás em Greenfoot?
\begin{choices}
 \choice `move(-5);'
 \choice `turn(180); move(5);'
 \choice `setLocation(getX(), getY() + 5);'
 \choice `moveBackwards(5);'
 \CC As alternativas `a' e `b' são funcionais.
\end{choices}

\question[1] Para alterar a imagem de um ator, qual método é usado?
\begin{choices}
 \choice `changeImage(new GreenfootImage("nova.png"));'
 \choice `setImage(new GreenfootImage("nova.png"));'
 \CC `setImage("nova.png");'
 \choice `updateImage("nova.png");'
 \choice As alternativas `b' e `c' são válidas.
\end{choices}

\question[1] Qual das seguintes afirmações sobre a classe `World' do Greenfoot é FALSA?
\begin{choices}
 \choice O construtor da classe `World' define suas dimensões.
 \choice Uma classe `World' pode ter seu próprio método `act()`.
 \choice A classe `World' pode adicionar e remover objetos do cenário.
 \CC Uma classe `World' herda da classe `Actor'.
 \choice O método `getObjects()' pode ser usado para encontrar atores no mundo.
\end{choices}


\question[1] Qual a principal desvantagem de usar `==' para comparar Strings?
\begin{choices}
 \choice É mais lento que `.equals()'.
 \choice Lança uma exceção se as strings forem diferentes.
 \CC Compara as referências de memória, não o conteúdo, o que pode levar a resultados inesperados.
 \choice Não funciona para strings criadas com `new String()'.
 \choice Só funciona se as strings estiverem em letras minúsculas.
\end{choices}

\question[1] Qual é a saída do código a seguir?
\lstinputlisting{code/Teste.java}
\begin{choices}
 \choice A
 \choice AB
 \choice AC
 \CC ACD
 \choice ABCD
\end{choices}

\question[1] Em Greenfoot, se um ator está na posição (10, 20) e executa `setLocation(getX() + 5, getY() - 5)', qual será sua nova posição?
\begin{choices}
 \choice (10, 20)
 \choice (5, 25)
 \CC (15, 15)
 \choice (5, 15)
 \choice (15, 25)
\end{choices}


\question[1] Qual das seguintes opções é a melhor prática para definir o nome de uma classe em Java?
\begin{choices}
 \choice `minhaClasse'
 \CC `MinhaClasse'
 \choice `minha\_classe'
 \choice `MINHA\_CLASSE'
 \choice `minha-classe'
\end{choices}

\question[1] Qual das seguintes opções é a melhor prática para o nome de uma constante em Java?
\begin{choices}
 \choice `minhaConstante'
 \choice `MinhaConstante'
 \choice `minha\_constante'
 \CC `MINHA\_CONSTANTE'
 \choice `minha-constante'
\end{choices}

\end{multicols}	
     \end{questions}	
\printmyanswers
\end{document}